\section{Incomplete areas \& future work}\label{sec:roadmap}

Every body section has reported its gaps in place; this section collects them into one ranked
inventory, records the work in flight and the branches parked behind it, separates what is
\emph{missing} from what was \emph{deliberately deferred}, and closes with a suggested order of
work --- which investments unlock which others.

\subsection{Incompleteness as a discipline}\label{sec:roadmap-signposting}

One finding from the incompleteness sweep deserves top billing: QtRocket's incompleteness is
unusually well-signposted. Almost every stub names the work that will consume it, and the roadmap
in \srcf{TODO.md} is accurate in both directions --- its completed-work claims verify against the
code, and its open items describe what is genuinely absent. The clearest example is the Barrowman
test suite, which says outright that it is pinning scaffolding for a consumer that does not exist
yet (\src{model/tests/AeroTests.cpp}{18}):

\begin{minted}{cpp}
// Exercises the Barrowman composition seam: per-part getAero folded by Part::getCompositeAero into a
// whole-rocket AeroProfile, with x_cp shared on the root-CM datum. NOTE: nothing in production
// consumes this in P2 -- these tests pin the seam for P5.
\end{minted}

\noindent The same discipline shows in the simulation core, where the future orientation
integrator sits as a commented-out \cppinline|std::unique_ptr<sim::RK4Solver<Quaternion>>| member
directly beside the live linear one, its interface already worked out in the accompanying comment:
it will reuse the \code{DESolver} abstraction with a time-keyed callback so
\cppinline|getTorques()| can be sampled at each stage's node time exactly as \code{getForces()} is
today (\src{sim/Propagator.h}{100}).

\begin{keyinsight}[Tested is not shipped]
The signposting cuts both ways. A contributor can locate the frontier by reading, because the
seams say what they are waiting for; but a reader skimming the tests could mistake the project for
further along than it is, because several finished artifacts have zero consumers --- the Barrowman
layer (unit-tested, uncalled; \src{model/tests/AeroTests.cpp}{18}), the wind model (compiled,
orphaned; \src{sim/WindModel.cpp}{16}), the torque interface (\src{model/RocketModel.cpp}{94}),
the quaternion state slots (\src{sim/StateData.h}{46}), and the ejection-delay data every motor
loader ingests (\src{model/RSEDatabaseLoader.cpp}{59}). The roadmap below is therefore mostly a
plan for \emph{consumption}, not construction.
\end{keyinsight}

\subsection{Consolidated inventory}\label{sec:roadmap-inventory}

Three tiers, ranked by user impact: what blocks a realistic flight, what blocks authoring and
everyday use, and what merely misleads a maintainer. Tier~1 collects the gaps that
\cref{sec:comparison} ranks first by flight-realism impact --- the machinery a realistic
powered flight cannot do without. Each entry is a summary with citations; the owning body section carries the analysis.

\subsubsection*{Tier 1 --- blocking realistic flight}

\begin{description}

  \item[Recovery phase and flight events --- \fabsent.] The only flight-phase logic in the
  codebase is the terminate predicate ``descending below the launch site''
  (\src{model/RocketModel.cpp}{65}); nothing changes at apogee, and no parachute or deployment
  identifier exists in production source. The data for ejection-triggered deployment is already
  ingested --- delays are parsed from RSE files (\src{model/RSEDatabaseLoader.cpp}{59}) and stored
  per motor (\src{model/MotorModel.h}{303}) --- but consumed by no flight logic; the
  thrustcurve.org client leaves them a TODO (\src{model/ThrustCurveClient.cpp}{198}). The typed
  termination taxonomy (\src{sim/Propagator.h}{35}) is the natural seed of an event vocabulary;
  see \cref{sec:sim} and \cref{sec:aero}.

  \item[Rotational dynamics --- \fabsent{} (dormant seams throughout).]
  \cppinline|getTorques()| returns zeros and has no production caller
  (\src{model/RocketModel.cpp}{94}); thrust is hardwired along world $+z$ regardless of attitude
  (\src{model/RocketModel.cpp}{71}); the orientation integrator is the commented-out member shown
  above (\src{sim/Propagator.h}{103}); and \code{StateData}'s quaternions are zero-initialized to
  a non-rotation, with a ``(vector, scalar)'' annotation contradicting Eigen's $(w,x,y,z)$
  constructor order (\src{sim/StateData.h}{46}). An angled launch tilts only the initial velocity.
  \Cref{sec:sim}, \cref{sec:aero}.

  \item[Barrowman seam unconsumed --- \fpartial{} (built, tested, uncalled).] Per-part
  \code{getAero} and the composite CP-on-the-CG-datum profile are complete and unit-tested, yet
  the only caller of \code{getCompositeAero} outside the model tests is itself a test
  (\src{tests/PlacementInvarianceTests.cpp}{89}). Flight drag remains one hand-entered line
  (\src{model/RocketModel.cpp}{88}), and \code{Propagatable::aeroData} is a default-constructed
  member nothing reads (\src{model/Propagatable.h}{61}). Two latent defects gate consumption: the
  \code{FinSet} chordwise sign mirror and the omitted child-tensor rotation (\cref{sec:aero},
  \cref{sec:model}).

  \item[Wind --- \fabsent{} (dead code).] \code{WindModel} is compiled into the sim library
  (\src{core/sim/CMakeLists.txt}{24}) but orphaned: it returns $(0,0,0)$ with all parameters ignored
  (\src{sim/WindModel.cpp}{16}), \code{Environment} has no wind member
  (\src{sim/Environment.h}{117}), and nothing instantiates it. Drag uses raw velocity, not
  relative airspeed (\src{model/RocketModel.cpp}{88}); \cref{sec:sim}.

  \item[Per-part drag build-up and reference area --- \fpartial.] Every concrete part's drag
  contribution is zero, so the per-part $C_d$ seam carries no signal (\cref{sec:aero});
  \code{deriveReferenceAreaFromGeometry()} correctly implements the max-frontal-disc convention
  but has no production caller (\src{model/parts/Part.cpp}{251}), so drag does not even scale with
  airframe class. No Mach or Reynolds number exists: the atmosphere's temperature and pressure are
  pinned by tests (\src{sim/tests/USStandardAtmosphereTests.cpp}{41}), but
  \code{USStandardAtmosphere::getSpeedOfSound} (\src{sim/USStandardAtmosphere.cpp}{143}) and its
  viscosity companion have no caller of any kind, production or test.

\end{description}

\subsubsection*{Tier 2 --- usability and authoring}

\begin{description}

  \item[GUI design editor --- \fpartial{} (display and persistence, no editor).] The merged
  \code{RocketTreeViewImplementation} work (\cref{sec:roadmap-inflight}) gives the main window a
  live read-only part tree bound through the Qt-free structure-changed callback
  (\src{gui/RocketTreeView.cpp}{149}) and working File Open/Save/Save~As
  (\src{gui/MainWindow.cpp}{95}); New and Close still ship \code{enabled=false}
  (\src{gui/MainWindow.ui}{118}), \code{RocketModelerView} is still an empty shell
  (\src{gui/RocketModelerView.cpp}{3}), and no part can be added, removed, or edited from any
  widget. Design authoring is CLI-only. \Cref{sec:interfaces}.

  \item[CLI placement authoring --- \fpresent{} (closed since the original review).] The
  \code{dynamic-Cd} merge gave \code{addpart} the full link vocabulary --- \code{seat=},
  \code{parentStation=}, \code{childStation=}, \code{gap=}, parent addressable by name
  (\src{cli/Repl.cpp}{221}) --- so every seat kind and station the model API and serializer
  support (\src{model/parts/Placement.h}{139}) is now authorable without hand-edited \code{.qrd}
  XML; the corpus scripts under \srcf{tests/data/designs/scripts/} exercise it end to end.
  \Cref{sec:interfaces}.

  \item[Motor placement, single motor, single stage --- \fpartial.] \code{motorOffset} is fixed
  at zero with no setter (\src{model/RocketModel.h}{148}), the motor's link is built from it
  (\src{model/RocketModel.cpp}{124}) and never persisted (\src{model/DesignSerializer.cpp}{94}),
  so every motor sits at the airframe CM, skewing CG$(t)$. One motor is structural
  (\code{findMotorInTree} is a first-match DFS, \src{model/RocketModel.cpp}{13}); one stage
  likewise --- \code{Stage.h} survives only as a ``Not yet'' commented include
  (\src{model/RocketModel.h}{23}), the original class removed in March 2024 (commit
  \code{46eca11}) to simplify prototyping. \Cref{sec:motors}.

  \item[Launch site --- \fabsent.] The spherical geoid ignores its latitude/longitude arguments
  (\src{sim/SphericalGeoidModel.cpp}{18}) and the gravity model caches ground level at a
  self-declared placeholder $(0,0)$ (\src{sim/SphericalGravityModel.cpp}{20}); no ECEF transform
  or launch-site input exists in any front end. \Cref{sec:sim}.

  \item[Simulation synchronous on the GUI thread --- \fabsent{} (no worker).]
  \code{launchRocket()} runs the whole propagation inline (\src{core/QtRocket.cpp}{54}) from a button
  slot (\src{gui/CannonballTab.cpp}{128}). Tolerable for a seconds-long ballistic hop; a
  minutes-long parachute descent --- the Tier~1 recovery work --- will freeze the UI, so that work
  forces this one. The broader GUI friction set is inventoried in \cref{sec:interfaces}.

\end{description}

\subsubsection*{Tier 3 --- hygiene and latent traps}

\begin{description}

  \item[Dead and dormant fields.] \code{Propagatable::nextState} is referenced nowhere
  (\src{model/Propagatable.h}{65}); \code{aeroData} is never read (\src{model/Propagatable.h}{61});
  \code{MotorModel::isp} is computed and dropped (\src{model/MotorModel.cpp}{130});
  \code{StateData}'s \code{dcm} and \code{eulerAngles} are never written
  (\src{sim/StateData.h}{49}); \code{MotorModel.h} retains commented-out serialization includes
  marked unused (\src{model/MotorModel.h}{10}). Each is a trap for a reader assuming it works.

  \item[Motor metadata losses.] The thrustcurve.org import leaves \code{certOrg} and delays
  unfilled with explicit TODOs (\src{model/ThrustCurveClient.cpp}{197}); \code{infoUrl} carries
  ``TODO: ???'' (\src{model/MotorModel.h}{307}); metadata members are public ``just for testing''
  (\src{model/MotorModel.h}{282}); typo'd certification strings are locked in for round-trip
  fidelity (\src{model/MotorModel.h}{144}); the RSE loader discards manufacturer-measured
  per-sample mass/CG data (\src{model/RSEDatabaseLoader.cpp}{89}). \Cref{sec:motors}.

  \item[Racy controller initialization.] The \code{QtRocket} singleton (\src{core/QtRocket.h}{29})
  guards construction with a mutex but tests its \code{initialized} flag \emph{before} taking the
  lock --- benign today only because every front end first-calls it from one thread; real the
  moment simulations move to a worker thread.

  \item[Documentation drift.] Four documents describe superseded designs: the persistence spec
  still shows the retired 0.1 \code{<offset>} schema
  (\src{docs/P2\_DESIGN\_PERSISTENCE\_AND\_CLI\_SPEC.md}{160}) that the shipped serializer rejects
  (\src{model/DesignSerializer.cpp}{176}); the visualizer's frozen spec documents the
  \code{Vector3}-offset scheme (\src{visualizer/SPEC.md}{38}) the visualizer itself no longer uses
  (\src{visualizer/RocketMesh.cpp}{434}); the architecture audit's gap register is half-stale
  (\src{docs/ARCHITECTURE\_AUDIT.md}{355}) with citations that no longer resolve
  (\src{sim/USStandardAtmosphere.h}{17}); \srcf{TODO.md} is accurate, but legacy phase numbers in
  code comments resolve only through its appendix translation table.

  \item[Internationalization and packaging vestiges.] \code{QTranslator} plumbing loads a
  locale-suffixed translation from a resource prefix that is never populated
  (\src{gui/main.cpp}{36}); distribution is a single template install rule with no packaging
  pipeline (\cref{sec:testing}).

  \item[Small tracked notes.] The remaining marked TODOs are honest and minor: \code{StateData}
  encapsulation (\src{sim/StateData.h}{29}), \code{utils::Bin} container-ification
  (\src{utils/Bin.h}{17}) and its linear-scan note (\src{utils/Bin.cpp}{39}).

\end{description}

\subsection{Work landed since the original review, and branch archaeology}\label{sec:roadmap-inflight}

This review is pinned to commit \code{20eca72}, the tip of \code{development}. The two branches
the original review tracked as work in flight have both landed. \code{dynamic-Cd} merged at
\code{fe11ae7}: the CLI seat vocabulary and \code{checkdesign}, the reworked
\code{DesignSerializer} \code{<link>} handling, the script-built fixture corpus
(\srcf{tests/data/designs/scripts/}), and the format-2 regeneration path for the
placement-invariance baseline (\cref{sec:persistence}). \code{RocketTreeViewImplementation}
merged at \code{29c9f9c} and was deleted: the live GUI part tree, the
\code{setStructureChangedCallback} seam, \code{Part::getParent()}, File Open/Save/Save~As, and
\code{CannonballTab::refreshFromModel()} (\cref{sec:interfaces}) --- the first slice of the
Tier~2 GUI editor, not the editor. The branch list has also been pruned since the original
review: the merged and abandoned experiments (\code{3Dviewer\_test},
\code{3Dviewer\_test\_codex}, \code{MotorPart}, \code{ConcretePartTypes}, \code{propagator},
\code{development-WIP}, \code{questionable\_changes}) are gone, leaving the short inventory of
\cref{tab:roadmap-branches}.

\begin{table}[htbp]
  \centering\small
  \begin{tabularx}{\linewidth}{@{}l c X@{}}
    \toprule
    \textbf{Remote branch} & \textbf{Ahead} & \textbf{Disposition} \\
    \midrule
    \code{development} & --- & Mainline; the reviewed commit \code{20eca72} is its tip. \\
    \code{dynamic-Cd} & $0$ & Merged at \code{fe11ae7} (seat vocabulary, corpus scripts); retained. \\
    \code{master} & behind & Retired old default branch. \\
    \code{website} & $+32$ & Jekyll project site; unrelated to the code. \\
    \bottomrule
  \end{tabularx}
  \caption{Remote branches at the review commit. ``Ahead'' counts commits not reachable from
  \code{development}.}
  \label{tab:roadmap-branches}
\end{table}

\subsection{Deliberately deferred, not forgotten}\label{sec:roadmap-deferred}

The Part Placement design document drew a distinction worth preserving: a deferred idea is a
planned extension point, not a dead end. Four of its deferrals are inherited by this roadmap, each
with its extension point already built.

\begin{description}

  \item[The joint graph.] Each child stores exactly one \code{StationLink} to one parent
  (\src{model/parts/Part.h}{292}), so a part bound to two hosts --- clusters, struts, pods --- is
  inexpressible. The implementation plan records this as ``deferred (correct long-term)''
  (\src{docs/part-placement-implementation-plan.md}{1031}): the single link per child can widen
  into a link list or joint table without disturbing the resolver's interface.

  \item[Off-axis seating.] A radial offset $r>0$ is used only for the seam and overlap
  \emph{check} while the pose solve keeps parts coaxial, because off-axis mass without an
  orientation would itself be a half-correct stored placement --- the error class the placement
  work eliminated. It arrives with 6-DOF, alongside \code{StationLink::childRot}, already carried
  as an identity quaternion (\src{model/parts/Placement.h}{60}), deliberately unpersisted
  (\src{model/DesignSerializer.cpp}{65}), and deliberately unapplied in the composite inertia walk
  --- the $R\,I\,R^{\mathsf{T}}$ rotation is omitted to stay bit-stable, with the code prescribing
  the test that must accompany its introduction (\src{model/parts/Part.cpp}{220}).

  \item[Fin sets with fewer than three fins.] \code{FinSet} warns and applies the $N \ge 3$
  axisymmetric approximation rather than the true anisotropic tensor
  (\src{model/parts/FinSet.cpp}{41}) --- honest under-modeling that becomes a real limit once
  rotational dynamics exist to feel the difference.

  \item[CLI auto-stack stations.] Shipped since the original review: \code{addpart}'s
  \code{parentStation=}/\code{childStation=} keys (\src{cli/Repl.cpp}{221}) express the stations
  the design plan once deferred, and the fixture scripts author seated, nested designs with them
  (\srcf{tests/data/designs/scripts/}). The deferral is retired.

\end{description}

\noindent These deferrals converge: the 6-DOF work consumes \code{childRot}, the tensor rotation,
and off-axis seating together, while the joint graph waits for a design that genuinely needs
multi-host attachment. Deferred is not dismissed --- each names the code already shaped to receive
it.

\subsection{Sequencing: what unlocks what}\label{sec:roadmap-sequencing}

\begin{designrule}[Consumption before construction]
Almost nothing in Tiers 1--2 requires new architecture. The aero profile, the torque interface,
the quaternion state, the ejection delays, the termination taxonomy, the authoring API, and the
serializer all exist and are tested. The highest-leverage sequencing wires existing seams to
consumers in dependency order, fixing the two known latent defects before their numbers become
user-visible.
\end{designrule}

\begin{enumerate}

  \item \textbf{Consume the Barrowman seam.} Prerequisites first: correct the \code{FinSet}
  chordwise sign convention, introduce the $R\,I\,R^{\mathsf{T}}$ rotation with its
  pinned-identity regression test (\src{model/parts/Part.cpp}{220}), and give
  \code{deriveReferenceAreaFromGeometry()} its missing caller (\src{model/parts/Part.cpp}{251}).
  The payoff is cheap and visible: a CP/CG/static-margin readout in CLI and GUI costs no new
  physics, exercises \code{getCompositeAero} for the first time, and flushes the sign defect while
  it is still a display bug rather than a flight bug (\cref{sec:aero}).

  \item \textbf{Recovery and flight events.} Independent of rotation, and the most visible
  fidelity gap: grow the termination taxonomy (\src{sim/Propagator.h}{35}) into an event
  vocabulary (liftoff, burnout, apogee, ejection, touchdown), consume the already-ingested
  ejection delays (\src{model/MotorModel.h}{303}), and switch drag configuration at deployment.
  This replaces the ballistic terminate at \src{model/RocketModel.cpp}{65} with a full flight ---
  and, because descents run minutes rather than seconds, it forces the worker-thread move off the
  GUI thread (\src{core/QtRocket.cpp}{54}), which in turn makes the initialization race worth fixing
  first.

  \item \textbf{Six degrees of freedom, then wind.} Reinstate the orientation integrator whose
  interface is already designed in place (\src{sim/Propagator.h}{103}), fixing the quaternion
  zero-initialization and the $(w,x,y,z)$-order ambiguity first (\src{sim/StateData.h}{46}); point
  thrust along the body axis instead of world $+z$ (\src{model/RocketModel.cpp}{71}); derive
  restoring torques from the profile consumed in step~1. Only then is wind worth plumbing: a point
  mass cannot weathercock, so connecting \code{WindModel} (\src{sim/WindModel.cpp}{16}) earlier
  would shift drag without producing any behavior wind exists to model.

  \item \textbf{The GUI design editor.} Pure consumption of a proven core: the CLI already
  demonstrates the full build--save--reload--fly loop against the same authoring API
  (\src{cli/Repl.cpp}{851}), the 25 tracked \code{.qrd} fixtures give the editor a ready test
  corpus, and the merged tree view already supplies the display surface and the
  change-notification bridge (\cref{sec:interfaces}). What remains is the property panel, part
  add/remove/edit, and the 2-D modeler --- the CLI's former seat/station gap is closed, so richer
  fixtures are already authorable for the GUI work itself.

\end{enumerate}

Staging and clusters (\cref{sec:motors}), the joint graph, and packaging stay off this critical
path deliberately: none blocks the four steps, and each gets cheaper after them --- staging needs
the event system, clusters need the joint graph, and the Tier~3 hygiene items are best retired by
whichever step touches their files. The pattern of the sequencing is the pattern of the codebase:
the hard structural machinery --- composite inertia, placement resolution, Barrowman composition,
adaptive integration, the termination taxonomy --- is built, tested, and waiting behind seams that
name their consumers. The rooms are framed and inspected; the remaining work is opening the doors
between them.
